\section{EOI（End of Interrupt）}

IRQ handler 执行完后\textbf{必须}向 PIC 发送 EOI，否则 PIC 认为当前中断
还在处理中，不会发送同级或更低优先级的后续中断。

\codefile{src/kernel/arch/pic.c}
\begin{lstlisting}[style=cstyle]
void pic_send_eoi(uint8_t irq) {
    if (irq >= 8) {
        outb(PIC2_CMD, 0x20);  // 先通知 slave
    }
    outb(PIC1_CMD, 0x20);      // 再通知 master（所有 IRQ 都要）
}
\end{lstlisting}

\begin{warnbox}
忘记发 EOI 是最常见的硬件中断 bug。表现为：第一次按键有反应，之后键盘完全死掉。
因为 PIC 还在等待上一次 IRQ1 的 EOI，拒绝发送新的中断。
\end{warnbox}

\begin{thinkbox}
为什么 slave IRQ（8--15）需要向\textbf{两片} PIC 都发 EOI？

因为 slave 是挂在 master 的 IRQ2 上的。一次 slave IRQ 同时占用了：
\begin{enumerate}
  \item slave PIC 的一个 IRQ 槽位
  \item master PIC 的 IRQ2 槽位
\end{enumerate}
两边都要通知释放。
\end{thinkbox}

% ============================================================================

